% Section 2: Model and Censoring Equivalence
\section{Model and Censoring Equivalence}
\label{sec:model}

\subsection{Model definition}
\label{sec:model-def}

Consider a parallel system of $m$ independent components with lifetimes
$T_1, \ldots, T_m$, where each component $j$ has an exponential distribution
with rate $\lambda_j > 0$:
\begin{equation}\label{eq:component-dist}
  T_j \sim \mathrm{Exp}(\lambda_j), \qquad
  F_j(t) = 1 - e^{-\lambda_j t}, \qquad
  f_j(t) = \lambda_j e^{-\lambda_j t}, \qquad t > 0.
\end{equation}
The system functions as long as at least one component survives, so the
system lifetime is
\begin{equation}\label{eq:system-lifetime}
  T = \max(T_1, \ldots, T_m).
\end{equation}
The system CDF and density are
\begin{align}
  F_T(t) &= \prod_{j=1}^m F_j(t) = \prod_{j=1}^m \bigl(1 - e^{-\lambda_j t}\bigr),
  \label{eq:system-cdf} \\
  f_T(t) &= \frac{d}{dt} F_T(t) = \sum_{j=1}^m f_j(t) \prod_{i \neq j} F_i(t),
  \label{eq:system-density}
\end{align}
the latter following from the product rule.  We write
$\blam = (\lambda_1, \ldots, \lambda_m) \in \Rset_{>0}^m$ for the parameter
vector.

Under \emph{Scheme~0} (black-box observation), we observe $n$ independent
copies $T^{(1)}, \ldots, T^{(n)}$ of the system lifetime~$T$.  No
component-level information is available: the observer sees only when the
system fails, not which components failed or in what order.  This is the
minimal observation scheme and the focus of the present paper.  More
informative schemes---periodic inspection, partial monitoring, continuous
component monitoring---are discussed briefly in \Cref{sec:discussion}.

\subsection{What system data reveals about components}
\label{sec:what-data-reveals}

When a system failure is observed at time $T = t$, the following facts
are logically entailed by the parallel structure~\eqref{eq:system-lifetime}:

\begin{enumerate}
  \item[(i)] \emph{All components have failed}: $T_j \leq t$ for every
    $j = 1, \ldots, m$.  Unlike the series system, where surviving
    components remain operational at system failure, the parallel system
    fails only after the last component expires.

  \item[(ii)] \emph{Exactly one component failed at time~$t$}: since the
    $T_j$ are independent and continuous, the probability of ties is zero,
    so almost surely there exists a unique index
    $J = \argmax_{j} T_j$ with $T_J = t$.

  \item[(iii)] \emph{The identity~$J$ is unobserved}: the system-level
    datum~$T = t$ does not reveal which component was the last to fail.
\end{enumerate}

\noindent
Each observation of $T$ carries the value of one component
lifetime exactly---but we do not know which one---and constrains all other
component lifetimes to lie below~$t$.  This structure is formalized in the
next subsection.

\subsection{The censoring equivalence}
\label{sec:censoring-equiv}

The following proposition recasts the incomplete-data structure of a parallel
system observation in the language of censoring.

\begin{proposition}[Censoring equivalence]\label{prop:censoring-equiv}
  The likelihood function of $\blam$ based on the system observation
  $T = \max(T_1, \ldots, T_m)$ is identical to the likelihood arising
  from the following component-level censoring scheme:
  \begin{enumerate}
    \item[\textup{(a)}] One component---the last to fail,
      $J = \argmax_j T_j$---is observed exactly: $T_J = T$.
    \item[\textup{(b)}] All other components are left-censored at~$T$: for
      $j \neq J$, we know only that $T_j < T$.
    \item[\textup{(c)}] The identity~$J$ is latent (unobserved).
  \end{enumerate}
\end{proposition}

\begin{proof}
Since $T = \max_j T_j$ and the component lifetimes are independent and
continuous, exactly one component satisfies $T_j = T$ almost surely; call
it~$J$.  For every other component, $T_j < T_J = T$.  The system-level
observation $T$ therefore carries precisely the following information:
\begin{itemize}
  \item the event $\{T_J = T\}$: one exact component observation;
  \item the events $\{T_j < T\}$ for all $j \neq J$: $m-1$ left-censored
    component observations at the common censoring time~$T$;
  \item the identity~$J$ is not observed.
\end{itemize}
No additional information is contained in the system observation $T = t$
beyond these three items, since $T = \max_j T_j$ is a deterministic function
of $(T_1, \ldots, T_m)$.  The complete data $(T_1, \ldots, T_m, J)$ augments
the observed data~$T$ with the latent identity and the left-censored lifetimes,
yielding the incomplete-data structure described in~(a)--(c).
\end{proof}

\Cref{prop:censoring-equiv} provides the conceptual foundation for the
estimation approach developed in \Cref{sec:likelihood,sec:estimation}.
The latent variable~$J$ plays a role analogous to the masked cause of
failure in series systems, but the censoring mechanism is fundamentally
different: left-censoring rather than right-censoring.
The result can also be obtained as a special case of the progressive
Type-II censoring representation for coherent systems
\citep{CramerNavarro2015}, which establishes that component failure times
from a system with known structure correspond to progressively censored
order statistics.  For a parallel system, this specializes to exactly
the one-exact plus $(m-1)$-left-censored structure of~(a)--(c).

\subsection{Contrast with series systems}
\label{sec:series-contrast}

\begin{remark}[Series--parallel duality in censoring]\label{rmk:series-parallel}
For a series system $T = \min(T_1, \ldots, T_m)$, the censoring equivalence
takes the dual form:
\begin{enumerate}
  \item[\textup{(a$'$)}] One component---the first to fail,
    $J = \argmin_j T_j$---is observed exactly: $T_J = T$.
  \item[\textup{(b$'$)}] All other components are \emph{right}-censored at~$T$:
    for $j \neq J$, we know only that $T_j > T$.
  \item[\textup{(c$'$)}] The identity~$J$ may be masked (observed, partially
    observed, or unobserved).
\end{enumerate}
The structural duality is: left-censoring in parallel systems corresponds to
right-censoring in series systems.  This mirrors the physical difference: in a
series system, surviving components have $T_j > T$ (they outlast the system); in
a parallel system, already-failed components have $T_j < T$ (they predecease the
system).

The notion of \emph{masking}---the conditions C1--C3 of
\citet{Miyakawa1984} and \citet{LinUsherGuess1993}---is physically motivated
for series systems, where the cause of system failure (which component failed
first) may be uncertain.  For parallel systems, masking is structurally
unmotivated: by the time the system fails, every component has failed, and
there is no single ``cause'' to mask.  Some authors
\citep{LiuShiZhangBai2017, RodriguesPereiraBheringPolpo2018} have
applied the C1--C3 framework to parallel systems, but the incomplete-data
problem is more naturally cast as component-level censoring with latent
identity, as in \Cref{prop:censoring-equiv}.  The censoring framing was
introduced by \citet{BheringPereiraPolpo2014} in a Bayesian Weibull setting;
we develop it here for frequentist exponential inference.
\end{remark}

\subsection{System-level censoring}
\label{sec:system-censoring}

In practice, the system lifetime~$T$ may itself be subject to censoring.
\Cref{prop:censoring-equiv} describes the component-level incomplete-data
structure \emph{given an exact observation of~$T$}.  When $T$ is censored,
the component-level implications compound as follows.

\begin{center}
\begin{tabular}{lll}
  \toprule
  \textbf{System observation} & \textbf{Component implication} & \textbf{Likelihood term} \\
  \midrule
  Exact: $T = t$ & $\exists!\, j\!: T_j = t$; $T_i < t\;\forall i \neq j$ &
    $f_T(t; \blam)$ \\[4pt]
  Right-censored: $T > \tau$ & $\exists\, j\!: T_j > \tau$; others unknown &
    $1 - F_T(\tau; \blam)$ \\[4pt]
  Interval: $a < T \leq b$ & All $T_j \leq b$; $\exists\, j\!: T_j > a$ &
    $F_T(b; \blam) - F_T(a; \blam)$ \\
  \bottomrule
\end{tabular}
\end{center}

\noindent
The right-censored case (system still functioning at monitoring time~$\tau$)
is the most common in practice.  Here, at least one component survives past
$\tau$, but we do not know which, nor do we know the status of the remaining
components: some may have failed silently before~$\tau$.  The system-level
survival probability $1 - F_T(\tau; \blam) = 1 - \prod_j (1 - e^{-\lambda_j \tau})$
has a closed form via the inclusion-exclusion expansion.

The interval-censored case arises under periodic inspection
(Scheme~1): the system fails between inspection times $a$ and $b$, so
$a < T \leq b$.  The likelihood contribution
$F_T(b; \blam) - F_T(a; \blam)$ is also available in closed form.

All three observation types---exact, right-censored, and
interval-censored---are accommodated by the likelihood framework of
\Cref{sec:likelihood}, where the relevant closed-form
expressions are developed in detail.
